\chapter{Introducción y contexto}
\label{chap:introduccion}

\section{Marco del proyecto}

Este Trabajo de Fin de Grado se desarrolla en el Grado en Ingeniería
Informática de la Facultat d'Informàtica de Barcelona, dentro de la especialidad
de Computación. El proyecto se sitúa en los ámbitos de la inteligencia
artificial y la visión por computador.

Su tema principal es el estudio de modelos de visión por computador
capaces de estimar las propiedades nutricionales a partir de imágenes de platos de comida, evaluando el error, la complejidad,
latencia y diversos factores.

El resultado del modelo debe interpretarse como una estimación experimental, no como un consejo médico.

\section{Motivación}

Registrar manualmente los alimentos y sus cantidades requiere tiempo y puede
introducir errores. Una revisión sistemática identificó las omisiones y la
estimación imprecisa del tamaño de las porciones entre las fuentes de error más
frecuentes del autorregistro \citep{whittonSystematicReviewExamining2022}. Un
estudio realizado en una muestra clínica concreta también observó desviaciones
relevantes; sus resultados no son generalizables, pero ilustran esta limitación
\citep{lichtmanDiscrepancySelfReportedActual1992}.

El uso de una imagen como entrada puede reducir parte de la carga del registro
manual y proporcionar una estimación complementaria. Sin embargo, esta
estimación debe evaluarse con rigor antes de utilizarse fuera de un contexto
experimental.

Este problema plantea una decisión técnica relevante. Los modelos de
mayor capacidad pueden obtener mejores predicciones, pero suelen necesitar más
memoria y tiempo de cálculo. En cambio, los modelos eficientes son más fáciles
de ejecutar de forma local, aunque pueden perder precisión. Estudiar este compromiso 
es también una de las motivaciones del proyecto.

\section{Contexto}

Los avances del aprendizaje profundo, la disponibilidad de conjuntos de datos
anotados y el desarrollo de modelos visuales han hecho posible investigar tareas
que relacionan el contenido de una imagen con las propiedades nutricionales de
un plato. Entre estas familias de modelos se encuentran las redes neuronales
convolucionales (CNN, por sus siglas en inglés) y los transformadores visuales.
En lugar de definir manualmente qué rasgos visuales son relevantes, estos
métodos aprenden patrones a partir de ejemplos anotados. Esto permite investigar
sistemas que reciben una fotografía como entrada y producen una estimación
nutricional, aunque la utilidad de esa estimación depende de sus errores, de los
datos disponibles y del contexto de uso. El estado del arte presenta los
conjuntos de datos y las arquitecturas concretas utilizados en la literatura.

El interés práctico reciente por registrar la ingesta a partir de fotografías
también se refleja en aplicaciones comerciales. MyFitnessPal ofrece la función
\emph{Meal Scan}, Healthify la función \emph{HealthifySnap} y Cal AI anuncia un
registro nutricional a partir de fotografías de comidas
\citep{myfitnesspalMealScan2026,healthifyHealthifySnap2026,calAiFoodCalorieTracker2026}.
Estas descripciones muestran la existencia de este tipo de funcionalidades,
pero no permiten cuantificar su adopción, no constituyen evidencia de precisión,
no permiten comparar las aplicaciones entre sí y no validan un uso clínico. El
TFG no evaluará dichas aplicaciones ni propondrá recomendaciones nutricionales
personalizadas.

\section{Identificación del problema}

El problema que se aborda es estimar a partir de una imagen RGB de un plato de comida, las propiedades
nutricionales de ese plato. Como propiedades nutricionales se consideran la masa total, las calorías y 
los principales macronutrientes (grasas, carbohidratos y proteínas).

Sin embargo, una fotografía tampoco muestra directamente toda la información necesaria. La cantidad de cada
ingrediente, la forma del plato o la presencia de ingredientes poco visibles
pueden producir errores importantes. Por tanto, el proyecto no parte de la
suposición de que un modelo visual eliminará los errores de estimación.

Otra dificultad presente es el sesgo de los datos de entrenamiento. Los datasets utilizados para entrenar
estos modelos de visión se centran en comida de una cultura o región concreta, restringiendose a un conjunto limitado de ingredientes y
tipos de platos concretos, por lo que el rendimiento puede variar significativamente dependiendo del plato que se evalue.

El resultado no solo se valorará con el error de predicción, sino también considerando el coste computacional,
la latencia, la complejidad del modelo, la arquitectura, el potencial de ejecución en un dispositivo con recursos
limitados y la relación entre estos factores.

\section{Conceptos básicos}

\subsection{Modelos eficientes y computación en el extremo}

La computación en el extremo, o \emph{edge computing}, consiste en ejecutar el
modelo cerca del lugar donde se generan los datos, por ejemplo en un teléfono o
en otro dispositivo local. Esta ejecución puede reducir la dependencia de una
conexión a Internet y evitar el envío continuo de imágenes. Sin embargo, obliga
a trabajar con límites de memoria, tiempo de respuesta y capacidad de cálculo.

\section{Actores implicados y usuarios potenciales}

En esta sección se distinguen los actores que pueden verse afectados y beneficiados 
por el proyecto:

\begin{description}
  \item[Comunidad académica.] Puede reutilizar el código, el protocolo y los
  resultados para estudiar modelos de estimación nutricional y ampliar el conocimiento en el área.
  \item[Servicios/Aplicaciones de registro alimentario.] Pueden integrar el modelo en sus sistemas,
  mejorando la precisión, el coste, o alguna otra métrica de interés. Traduciendose en una mejor experiencia
  o abaratamiento de costes para ofrecer esta funcionalidad.
  \item[Usuarios potenciales.] Personas que quieran utilizar sistemas de registro alimentario
  para llevar un control de su ingesta por temas deportivos y de salud.
  \item[Proveedores de datos.] Mantienen los conjuntos de datos de los que
  dependen el entrenamiento y la evaluación.
\end{description}
